\section{Dekooderipohjainen Transformer-arkkitehtuuri suurille
kielimalleille}\label{sec:transformer}

Luvuissa \ref{sec:discrete} ja \ref{sec:mlp} kielimallinnusta
tarkasteltiin seuraavan tokenin ehdollisen todennäköisyyden estimointina.
Olkoon \(\mathcal{V}\) tokenisanasto ja
\(\Delta(\mathcal{V})\) sanaston yli määriteltyjen
todennäköisyysjakaumien joukko. Kielimalli voidaan tällöin kuvata
funktiona, joka muuntaa aiemmin havaitun tokenikontekstin seuraavan
tokenin todennäköisyysjakaumaksi.

\(n\)-gram-malli ennustaa tokenin \(n-1\) mittaisen kontekstin
perusteella. Se voidaan siten kuvata funktiona
\begin{align*}
  f_{\mathrm{ngram}}
  \colon
  \mathcal{V}^{n-1}
  \to
  \Delta(\mathcal{V}).
\end{align*}

Luvussa \ref{sec:mlp} tarkastellussa MLP-kielimallissa syötekerroksen dimensio sitoo mallin tiettyyn kontekstin pituuteen. Transformer-mallit perustuvat samaan ennustustehtävään, mutta tokenien esityksiä käsitellään sekvenssinä, jonka eri positioihin sovelletaan samoja opittavia muunnoksia. Sama parametrisaatio soveltuu näin eripituisiin syötteisiin mallin konteksti-ikkunan rajoissa.

Olkoon \(L\) mallin suurin sallittu
kontekstinpituus. Oletetaan, että mallin syöte on aina epätyhjä tokenijono (tyhjä syöte korvataan aloitustokenilla $\langle s \rangle$), jolloin dekooderipohjainen Transformer-kielimalli voidaan
kuvata funktiona
\begin{align*}
  f_{\mathrm{Tr}}
  \colon
  \bigcup_{m=1}^{L}\mathcal{V}^{m}
  \to
  \Delta(\mathcal{V}).
\end{align*}
Malli voi siten muodostaa seuraavan tokenin jakauman millä tahansa
enintään \(L\) tokenin pituisella kontekstilla. Käytännössä suurin
kontekstinpituus on edelleen arkkitehtuurin ja laskentaresurssien
rajoittama, mutta mallin parametrien ei tarvitse olla sidottuja yhteen
tiettyyn kontekstipituuteen.

Tässä luvussa tarkastellaan erityisesti dekooderipohjaista
Transformer-arkkiteh-\\tuuria, joka muodostaa useiden generatiivisten
kielimallien, kuten GPT-malliper-\\heen, matemaattisen perustan \cite{radford2018improving}.
Transformeriin perustuvan arkkitehtuurin esitteli alun perin Vaswanin ym. vuonna 2017
\cite{vaswani2017attention}. Alkuperäinen arkkitehtuuri oli
konekäännökseen suunniteltu enkooderi--dekooderi-rakenne. Pian tämän jälkeen esitettiin, että pelkkään Transformer-dekooderiin perustuva autoregressiivinen kielimalli voidaan esikouluttaa suurilla tekstiaineistoilla ja hienosäätää useisiin kielenkäsittelytehtäviin \cite{radford2018improving}.

\subsection{Itsehuomio}

Transformer-mallin keskeinen mekanismi on itsehuomio (engl.\ \emph{self-attention}). Sen avulla tokenien merkitys ei määräydy pelkän etäisyyden perusteella, vaan malli painottaa dynaamisesti kunkin tokenin kannalta olennaista tietoa koko sekvenssin laajuudelta, kun se päivittää upotusvektoreita.

Olkoon tarkasteltavan tokenijonon pituus \(T\) ja olkoon \(h_i \in \mathbb{R}^{1 \times d}\) position \(i\) tokenia vastaava esitysvektori. Näistä muodostetaan syötematriisi
\begin{align*}
  H
  =
  \begin{bmatrix}
    h_1 \\
    h_2 \\
    \vdots \\
    h_T
  \end{bmatrix}
  \in \mathbb{R}^{T \times d}.
\end{align*}

Itsehuomiomekanismi muodostaa syöte-esityksistä kolme erillistä lineaarista projektiota: kyselyn (engl.\ \emph{query}), avaimen (engl.\ \emph{key}) ja arvon (engl.\ \emph{value}). Olkoot \(W_Q, W_K \in \mathbb{R}^{d \times d_k}\) ja \(W_V \in \mathbb{R}^{d \times d_v}\) opittavia painomatriiseja. Kysely-, avain- ja arvomatriisit lasketaan kertomalla syötematriisi näillä painoilla:
\begin{align*}
  Q &= H W_Q \in \mathbb{R}^{T \times d_k}, \\
  K &= H W_K \in \mathbb{R}^{T \times d_k}, \\
  V &= H W_V \in \mathbb{R}^{T \times d_v}.
\end{align*}
Kyselymatriisi määrittää haettavan informaation tyypin kunkin tokenin kohdalla, kun taas avainmatriisi kuvaa tokenin tarjoamaa kontekstitietoa. Arvomatriisi puolestaan edustaa varsinaista informaatiosisältöä, josta uudet kontekstualisoidut esitykset muodostetaan huomiopainotusten avulla.

Kyselyiden ja avainten välinen vastaavuus lasketaan pistetuloina, joista muodostuu skaalattu huomiopistematriisi
\begin{align*}
  S = \frac{QK^{\top}}{\sqrt{d_k}} \in \mathbb{R}^{T \times T}.
\end{align*}
Matriisin alkio \(S_{ij}\) mittaa position \(i\) kyselyn ja
position \(j\) avaimen yhteensopivuutta. Skaalausta voidaan
perustella tarkastelemalla kysely- ja avainvektoreita, joiden
komponentit ovat riippumattomia, nollakeskiarvoisia ja
yksikkövarianssisia. Näillä oletuksilla niiden pistetulon
varianssi on \(d_k\), joten jakaminen luvulla \(\sqrt{d_k}\)
palauttaa varianssin arvoon yksi. Vaikka oletukset eivät
yleisesti toteudu täsmällisesti opituilla esityksillä,
skaalaus hillitsee huomiopisteiden suuruuden kasvua dimension
mukana. Tämä vähentää softmax-funktion saturoitumisen riskiä
ja tukee vakaampaa gradienttipohjaista koulutusta
\cite{vaswani2017attention}.

Huomiopisteet normalisoidaan todennäköisyysjakaumiksi soveltamalla softmax-funktiota riveittäin, jolloin saadaan huomiopainomatriisi
\begin{align*}
  A = \operatorname{softmax}(S) = \operatorname{softmax}\left(\frac{QK^{\top}}{\sqrt{d_k}}\right) \in \mathbb{R}^{T \times T},
\end{align*}
missä kukin alkio \(A_{ij} \in [0, 1]\) ilmaisee, kuinka suuri painoarvo position \(j\) arvolle annetaan position \(i\) uutta esitystä laskettaessa, ja riveille pätee \(\sum_{j=1}^{T} A_{ij} = 1\).

Lopullinen itsehuomion tulos lasketaan arvomatriisin rivien
painotettuna summana matriisitulon avulla
\begin{align*}
  \operatorname{Attention}(H)
  =
  AV
  =
  \operatorname{softmax}
  \left(
    \frac{QK^{\top}}{\sqrt{d_k}}
  \right)V
  \in
  \mathbb{R}^{T \times d_v}.
\end{align*}
Tuloksena saadun matriisin kukin rivi \(i\) edustaa position \(i\)
uutta kontekstualisoitua esitystä, johon on koottu tietoa kaikista
sekvenssin tokeneista huomiopainojen \(A_{ij}\) mukaisesti.

Autoregressiivisissä kielimalleissa tavoitteena on mallintaa seuraavan tokenin ehdollista todennäköisyyttä, joten tokenin esitys ei voi nojata tulevaan kontekstiin. Tätä varten itsehuomiota rajoitetaan kausaalisella maskilla (engl. \emph{causal mask}), joka estää informaation vuotamisen tulevaisuudesta asettamalla position \(i\) huomiopainot nollaksi kaikille positioille \(j > i\). Tämä toteutetaan lisäämällä huomiopisteisiin ennen
softmax-funktiota maskimatriisi
\(M \in (\mathbb{R} \cup \{-\infty\})^{T \times T}\),
jonka alkiot määritellään muodossa
$$
M_{ij} =
\begin{cases}
  0, & \text{kun } j \le i, \\
  -\infty, & \text{kun } j > i.
\end{cases}
$$
Kausaalinen itsehuomio lasketaan tällöin muodossa
\begin{align*}
  \operatorname{CausalAttention}(Q,K,V) = \operatorname{softmax}\left(\frac{QK^{\top}}{\sqrt{d_k}} + M\right)V.
\end{align*}
Käytännöllä \(\exp(-\infty) = 0\) huomiopainomatriisista muodostuu alakolmiomatriisi. Tällöin kunkin tokenin esitys riippuu vain siitä itsestään ja sitä edeltävistä tokeneista, mikä säilyttää generointiprosessin kausaalisen rakenteen.

Yksi huomiopää kokoaa informaatiota yhden opitun kysely-, avain- ja
arvoprojektion varassa, ja sama huomiopainotus \(A_{ij}\) koskee
arvovektorien kaikkia koordinaatteja. Monipäinen itsehuomio
(engl.\ \emph{multi-head attention}) mahdollistaa useiden
erilaisten painotusten ja esitysten muodostamisen rinnakkain \cite{vaswani2017attention}.

Olkoon päiden lukumäärä \(n_\text{heads}\). Kukin pää
\(k \in \{1,\dots,n_\text{heads}\}\) suorittaa kausaalisen itsehuomion
omilla opittavilla projektioillaan
$$
\operatorname{head}_k
=
\operatorname{CausalAttention}
\left(
  H W_k^Q,\,
  H W_k^K,\,
  H W_k^V
\right),
$$
missä \(W_k^Q, W_k^K \in \mathbb{R}^{d \times d_k}\) ja
\(W_k^V \in \mathbb{R}^{d \times d_v}\). Kukin pää muodostaa näin oman huomiopainomatriisinsa ja omat painotetut summansa arvovektoreista. Päät voivat oppia toisiaan täydentäviä riippuvuuksia. Esimerkiksi yksi pää saattaa keskittyä syntaktisiin rakenteisiin ja toinen laajempiin semanttisiin yhteyksiin. Todellisuudessa opitut suhteet eivät kuitenkaan jakaudu näin siististi, vaan ne ovat usein huomattavasti monimutkaisempia ja vaikeammin tulkittavia \cite{clark2019what}.

Lopuksi eri päiden tuottamat esitykset ketjutetaan ja
yhdistetään lineaarisella projektiolla takaisin mallin
ulottuvuuteen \(d\):
$$
\operatorname{CausalMultiHead}(H)
=
\left[
  \operatorname{head}_1
  \,\|\,
  \operatorname{head}_2
  \,\|\,
  \dots
  \,\|\,
  \operatorname{head}_{n_\text{heads}}
\right] W^O,
$$
missä \(\|\) merkitsee sarakkeittaista ketjuttamista ja
\(W^O \in \mathbb{R}^{n_\text{heads} d_v \times d}\) on opittava
parametrimatriisi.

Monipäisyyden etu on että eri päät
voivat käyttää toisistaan poikkeavia huomiopainotuksia ja
arvoprojektioita, joiden tuottama informaatio yhdistetään
tulosprojektiossa. Käytännössä pään dimensioiksi valitaan
usein \(d_k = d_v = d/n_\text{heads}\), jolloin monipäisen itsehuomion
laskennallinen vaativuus vastaa karkeasti täyden dimension
yksipäistä huomiota \cite{vaswani2017attention}.

\subsection{Positiokoodaus}

Maskiton itsehuomio on permutaatioekvivariantti eli ilman erillistä
positiotietoa se ei kykene erottelemaan tokenien järjestystä. Olkoon
\(P \in \mathbb{R}^{T \times T}\) permutaatiomatriisi, joka järjestää
syötematriisin \(H\) rivit uudelleen. Koska kysely-, avain- ja
arvomatriisit muodostetaan riveittäin samoilla lineaarisilla
kuvauksilla, permutoidusta syötteestä saadaan
\(Q' = PQ\), \(K' = PK\) ja \(V' = PV\). Tällöin
huomiopistematriisille pätee
\begin{align*}
  \frac{Q'(K')^{\top}}{\sqrt{d_k}}
  =
  P
  \frac{QK^{\top}}{\sqrt{d_k}}
  P^{\top}.
\end{align*}

Riveittäin sovellettava softmax-funktio säilyttää vastaavan
permutaation, joten maskittomalle itsehuomiolle pätee
\begin{align*}
  \operatorname{Attention}(PH)
  =
  P\operatorname{Attention}(H).
\end{align*}

Syötteen tokenien permutointi siis ainoastaan permutoi tulosteet
samalla tavalla. Maskiton itsehuomiomekanismi ei sellaisenaan sisällä
tietoa siitä, missä järjestyksessä tokenit esiintyvät.
Autoregressiivisessa mallissa kausaalinen maski rikkoo tämän
ominaisuuden mielivaltaisten permutaatioiden suhteen, sillä se määrää,
mitkä positiot ovat kunkin tokenin menneisyydessä. Kausaalinen maski
tuo siten malliin järjestysrakennetta, jonka avulla malli voi myös
oppia hyödyntämään sijaintiin liittyvää tietoa. Maski ei kuitenkaan
sisällytä kyselyn ja avaimen väliseen yhteensopivuuteen eksplisiittistä
tietoa siitä, sijaitseeko avain esimerkiksi yhden vai usean tokenin
päässä kyselystä. Erillinen positiokoodaus antaa mallille suoran
mekanismin tokenien sijaintien tai niiden välisten positioerojen
esittämiseen.

Positiotieto voidaan sisällyttää Transformer-arkkitehtuuriin usealla
tavalla. Alkuperäisessä Transformerissa tokeniesityksiin lisätään
absoluuttiset, sinimuotoiset positiokoodaukset
\cite{vaswani2017attention}. Suhteellisissa menetelmissä
huomiopisteitä puolestaan muokataan positioiden välisen etäisyyden
perusteella. Monissa moderneissa dekooderipohjaisissa kielimalleissa
käytetään kiertävää positiokoodausta (engl.\ \emph{rotary positional
embedding}, RoPE), jossa positiotieto liitetään kysely- ja
avainvektoreihin rotaatioiden avulla \cite{su2024roformer}.

Olkoon yhden huomiopään kysely- ja avainvektorien dimensio \(d_k\)
parillinen. RoPE jakaa vektorien koordinaatit \(d_k/2\)
kaksiulotteiseen pariin ja kiertää jokaista paria position mukaisella
kulmalla. Kun parit indeksoidaan arvoilla
\(\ell \in \{0,\dots,d_k/2-1\}\), tavanomainen kulmataajuus pareille on
\begin{align*}
  \omega_{\ell}
  =
  10000^{-2\ell/d_k}.
\end{align*}

Kulmaa \(\phi\) vastaava kaksiulotteinen rotaatiomatriisi on

$$
R(\phi)
=
\begin{bmatrix}
  \cos \phi & -\sin \phi \\
  \sin \phi & \cos \phi
\end{bmatrix}.
$$

Position \(i\) koko rotaatiomatriisi muodostetaan näistä
kaksiulotteisista rotaatioista lohkodiagonaalisena matriisina

$$
R_i
=
\operatorname{diag}
\left(
  R(i\omega_0),
  R(i\omega_1),
  \dots,
  R\left(i\omega_{d_k/2-1}\right)
\right)
\in
\mathbb{R}^{d_k \times d_k}.
$$

Olkoon \(q_i = Q_{i,:} \in \mathbb{R}^{1 \times d_k}\) position \(i\) kyselyvektori ja
\(k_j = K_{j,:} \in \mathbb{R}^{1 \times d_k}\) position \(j\) avainvektori. RoPE korvaa
nämä kierretyillä rivivektoreilla
$$
\tilde q_i = q_i R_i,
\qquad
\tilde k_j = k_j R_j.
$$

Arvovektoreihin rotaatiota ei tavallisesti sovelleta. Kierrettyjen
kysely- ja avainvektorien pistetuloksi saadaan
$$
\tilde q_i \tilde k_j^{\top}
=
q_i R_i (k_j R_j)^{\top}
=
q_i R_i R_j^{\top} k_j^{\top}.
$$

Kaksiulotteisille rotaatiomatriiseille pätee
\(R(\alpha)R(\beta)^{\top}=R(\alpha-\beta)\). Sama ominaisuus pätee
lohkoittain myös RoPE-rotaatioille, joten
$$
R_i R_j^{\top}
=
R_{i-j}.
$$

Siten RoPE-huomiopiste voidaan kirjoittaa muodossa
$$
\tilde q_i \tilde k_j^{\top}
=
q_i R_{i-j} k_j^{\top}.
$$

Vaikka kysely ja avain kierretään niiden absoluuttisten positioiden
perusteella, rotaatioiden yhteisvaikutus pistetulossa riippuu ainoastaan
positioiden suhteellisesta siirtymästä \(i-j\).

Kun kierretyt kyselyt ja avaimet kootaan matriiseiksi
\(\tilde Q\) ja \(\tilde K\), kausaalinen itsehuomio saa muodon

$$
\operatorname{CausalAttention}
\left(
  \tilde Q,\tilde K,V
\right)
=
\operatorname{softmax}
\left(
  \frac{\tilde Q\tilde K^{\top}}{\sqrt{d_k}}
  +
  M
\right)V.
$$

Koska \(R_i\) on ortogonaalinen matriisi, rotaatio säilyttää kysely- ja avainvektorien normit. Eri koordinaattipareille määritellyt kulmataajuudet määräävät, kuinka nopeasti rotaatiokulma muuttuu position funktiona. Suurilla taajuuksilla kulma muuttuu nopeasti position kasvaessa, kun taas pienillä taajuuksilla muutos on hitaampaa. Näin eri koordinaattiparit kuvaavat positioeroja eri mittakaavoilla. RoPE ei muuta itsehuomion asymptoottista laskennallista vaativuutta \cite{su2024roformer}.

\subsection{Dekooderikerros}

Yksittäinen dekooderikerros koostuu kausaalisesta monipäisestä
itsehuomiosta ja positiokohtaisesta myötäkytkentäverkosta (engl.\ \emph{feed-forward network}, FFN).
Kumpaankin komponenttiin liitetään jäännöskytkentä
(engl.\ \emph{residual connection}) sekä normalisointi.
Normalisointi auttaa hallitsemaan esitysten mittakaavaa,
ja jäännöskytkennät helpottavat gradienttien kulkua syvässä
verkossa. Yhdessä nämä rakenteet tukevat koulutuksen
vakautta \cite{he2016deep, ba2016layer, vaswani2017attention}.

Laskennan vakauttamiseksi esitykset normalisoidaan ennen kuhunkin
alilohkoon syöttämistä. Normalisointifunktio \(\operatorname{Norm} \colon
\mathbb{R}^{1 \times d} \to \mathbb{R}^{1 \times d}\) sovelletaan itsenäisesti kunkin position
esitysvektoriin. Yleinen valinta nykyaikaisissa malleissa on neliöllisen
keskiarvon normalisointi (engl.\ \emph{root mean square normalization},
RMSNorm) \cite{zhang2019root}. Tämä skaalaa vektorin sen komponenttien neliöllisen keskiarvon avulla
$$
\operatorname{RMSNorm}(x) = \frac{x}{\operatorname{RMS}(x)} \odot \gamma,
$$
missä \(\odot\) on Hadamard-tulo, \(\gamma \in \mathbb{R}^{1 \times d}\)
on opittava skaalausparametri ja
$$
\operatorname{RMS}(x) = \sqrt{\frac{1}{d} \sum_{j=1}^{d} x_j^2 + \epsilon},
$$
missä \(\epsilon > 0\) estää nollalla jakamisen. Matriisille
\(H \in \mathbb{R}^{T \times d}\) normalisointi \(\operatorname{Norm}(H)\)
määritellään soveltamalla kuvausta riveittäin jokaiseen esitysvektoriin.

Itsehuomiokerrosta seuraa positiokohtainen myötäkytkentäverkko. Toisin kuin luvussa \ref{sec:mlp}
tarkastellussa MLP-kielimallissa, jossa verkko muuntaa koko kontekstin suoraan
seuraavan tokenin logiteiksi, Transformer-lohkossa FFN toimii positiokohtaisena
muunnoksena. Sitä sovelletaan itsenäisesti ja samoilla parametreilla kunkin
position esitykseen erikseen. Positioiden välisestä informaation siirrosta
vastaa siten itsehuomio, kun taas FFN vastaa yksittäisen esityksen rikastamisesta
ja epälineaarisesta muuntamisesta.

Kaksi affiinia muunnosta sisältävä FFN voidaan määritellä
vektorille \(x \in \mathbb{R}^{1 \times d}\) muodossa
$$
\operatorname{FFN}(x)
=
\phi(x W_1 + b_1) W_2 + b_2,
$$
missä \(W_1 \in \mathbb{R}^{d \times d_{ff}}\),
\(W_2 \in \mathbb{R}^{d_{ff} \times d}\),
\(b_1 \in \mathbb{R}^{1 \times d_{ff}}\),
\(b_2 \in \mathbb{R}^{1 \times d}\) ja \(\phi\) on
koordinaateittain sovellettava epälineaarinen aktivaatiofunktio.
Sisäisen piiloesityksen dimensio \(d_{ff}\) valitaan
tyypillisesti suuremmaksi kuin mallin dimensio \(d\).
Alkuperäisessä Transformer-arkkitehtuurissa käytettiin
valintaa \(d_{ff} = 4d\) \cite{vaswani2017attention}.

Kun tokenien esitykset ovat matriisin
\(H \in \mathbb{R}^{T \times d}\) riveinä, sama laskenta
kirjoitetaan muodossa
$$
\operatorname{FFN}(H)
=
\phi\left(
  H W_1 + \mathbf{1} b_1
\right) W_2
+
\mathbf{1} b_2,
$$
missä \(\mathbf{1} \in \mathbb{R}^{T \times 1}\) on sarakeykkösvektori,
jonka avulla bias-vektorit lisätään jokaiselle riville.

Koko dekooderilohko voidaan määritellä funktiona
\(\operatorname{DecoderBlock} \colon \mathbb{R}^{T \times d} \to \mathbb{R}^{T \times d}\).
Olkoon syöte \(H \in \mathbb{R}^{T \times d}\). Lohkon laskenta
etenee kahdessa vaiheessa:
$$
H^{(1)} = H + \operatorname{CausalMultiHead}(\operatorname{Norm}_1(H)),
$$
$$
\operatorname{DecoderBlock}(H) = H^{(1)} + \operatorname{FFN}(\operatorname{Norm}_2(H^{(1)})),
$$
missä \(\operatorname{Norm}_1\) ja
\(\operatorname{Norm}_2\) ovat toisistaan erillisiä
normalisointeja omilla opittavilla parametreillaan.
Normalisointi sijoitetaan ennen kumpaakin alilohkoa, joten rakenne
noudattaa niin kutsuttua \emph{pre-norm}-järjestystä
\cite{xiong2020layer}.

Jäännöskytkentä muodostaa gradientille suoran kulkureitin
alilohkon ohi, ja normalisointi estää gradienttien räjähtämisen tai vaimenemisen kerrosten kasvaessa, mikä helpottaa syvien verkkojen optimointia
\cite{he2016deep} \cite{xiong2020layer}.

\subsection{Dekooderikerroksista seuraavan tokenin jakaumaan}

Yksittäinen dekooderilohko muuntaa tokenien esitysmatriisin toiseksi samankokoiseksi esitysmatriisiksi. Kokonaisessa kielimallissa tokenien upotuksesta saatua esitysmatriisia käsitellään peräkkäisissä dekooderilohkoissa, minkä jälkeen tuloskerros muuntaa lopulliset esitykset seuraavan tokenin todennäköisyysjakaumaksi seuraavalle tokenille.

Olkoon syötteenä havaittu konteksti \(w_{1:T} = (w_1, \dots, w_T)\),
missä \(1 \le T \le L\). Kuten luvussa \ref{sec:mlp} määriteltiin, kukin
token muunnetaan esitysvektoriksi upotusmatriisin
\(E \in \mathbb{R}^{|\mathcal{V}| \times d}\) avulla. Ensimmäisen lohkon
syötematriisi \(H^{[0]} \in \mathbb{R}^{T \times d}\) muodostetaan näistä
riveittäin
$$
H^{[0]}_{i,:} = e(w_i) = x(w_i)E, \qquad i = 1, \dots, T.
$$

Tässä tarkasteltavassa arkkitehtuurissa positiotieto sisällytetään
itsehuomion ky\\-sely- ja avainvektoreihin RoPE-rotaatioiden avulla.
Jatkossa dekooderilohkon kausaalisen monipäisen
itsehuomion oletetaan sisältävän edellä määritellyt rotaatiot
kunkin huomiopään kyselyille ja avaimille.

Olkoon dekooderilohkojen lukumäärä \(B\). Esitysmatriisia
muunnetaan peräkkäisissä lohkoissa seuraavasti
$$
H^{[b]}
=
\operatorname{DecoderBlock}_b
\left(H^{[b-1]}\right),
\qquad
b=1,\dots,B.
$$
Kullakin lohkolla on omat opittavat parametrinsa, mutta
esitysmatriisin koko säilyy kaikissa vaiheissa samana.
Hakasulkeinen yläindeksi ilmaisee tässä koko mallin
dekooderilohkon indeksin, eikä yksittäisen lohkon sisäistä
laskentavaihetta.

Kausaalisuus säilyy myös useita lohkoja yhdistettäessä.
Ensimmäisessä lohkossa position \(i\) esitys voi hyödyntää
ainoastaan positioita \(1,\dots,i\). Seuraavassa lohkossa
itsehuomio yhdistää vain näihin positioihin kuuluvia esityksiä,
joista yksikään ei sisällä tietoa positiota \(i\) seuraavista
tokeneista. Koska normalisointi ja myötäkytkentäverkko toimivat
positiokohtaisesti, ne eivät muuta tätä riippuvuusrakennetta.
Siten position \(i\) esitys riippuu kaikissa lohkoissa vain
kontekstista \(w_{1:i}\).

Tarkasteltavassa pre-norm-rakenteessa viimeisen dekooderilohkon
tulokseen sovelletaan positiokohtaista normalisointia
$$
Z
=
\operatorname{Norm}_{\mathrm{out}}
\left(H^{[B]}\right)
\in
\mathbb{R}^{T \times d}.
$$
Olkoon \(z_i = Z_{i,:} \in \mathbb{R}^{1 \times d}\) position \(i\)
lopullinen esitys vektorina. Se muunnetaan sanaston
ulottuvuuteen affiinilla kuvauksella
$$
\ell_i
=
z_i W_{\mathrm{out}} + b_{\mathrm{out}},
$$
missä \(W_{\mathrm{out}} \in \mathbb{R}^{d \times |\mathcal{V}|}\)
ja \(b_{\mathrm{out}} \in \mathbb{R}^{1 \times |\mathcal{V}|}\)
ovat opittavia parametreja. Vektorin \(\ell_i\) komponentti
\(\ell_{i,j}\) vastaa sanaston tokenin \(v_j\) normalisoimatonta
pistemäärää eli logittia. Näistä muodostetaan todennäköisyysjakauma
softmax-funktiolla.

Merkitään mallin kaikkia opittavia parametreja symbolilla \(\theta\).
Position \(i\) piiloesitys tiivistää informaation kontekstista \(w_{1:i}\),
joten sen tuottama logittivektori \(\ell_i\) määrittää todennäköisyysjakauman
seuraavalle tokenille \(w_{i+1}\)
$$
p_{\theta}(w_{i+1}=v_j \mid w_{1:i})
=
\frac{\exp(\ell_{i,j})}
{\sum_{k=1}^{|\mathcal{V}|}\exp(\ell_{i,k})},
\qquad
v_j \in \mathcal{V}.
$$
Koko syötekontekstia \(w_{1:T}\) seuraavan uuden tokenin \(w_{T+1}\) jakauma
saadaan siten suoraan viimeisen position logiteista
$$
f_{\mathrm{Tr}}(w_{1:T})
=
\operatorname{softmax}(\ell_T)
\in
\Delta(\mathcal{V}).
$$
Kielimallin laskentaketju etenee näin diskreeteistä tokeneista
jatkuviin esityksiin, kontekstuaalisiin muunnoksiin ja lopulta
takaisin diskreetin sanaston yli määritellyksi jakaumaksi.

Arkkitehtuuri määrittää seuraavan tokenin jakauman annetuilla
parametreilla \(\theta\). Seuraavaksi tarkastellaan, miten
parametrit opitaan aineistosta ja miten opittua jakaumaa
käytetään kokonaisten tokenijonojen muodostamiseen.

\subsection{Autoregressiivinen koulutus ja generointi}

Transformer-kielimallin koulutus perustuu samaan
uskottavuusperiaatteeseen kuin luvussa \ref{sec:mlp} tarkastellussa MLP-kielimallissa. Todennäköisyyksien ketjusäännön perusteella tokenijonon
\(w_{1:T}\) todennäköisyys voidaan esittää muodossa
$$
p_{\theta}(w_{1:T})
=
\prod_{t=1}^{T}
p_{\theta}(w_t \mid w_{<t}),
$$
missä \(w_{<t}=(w_1,\dots,w_{t-1})\).

Olkoon koulutusaineisto
\(\mathcal{D}=\{w^{(1)},\dots,w^{(N)}\}\), missä
\(w^{(r)}=(w_1^{(r)},\dots,w_{T_r}^{(r)})\) ja
\(1 \le T_r \le L\). Aineiston tokenikohtainen
keskimääräinen negatiivinen log-uskottavuus on
$$
J(\theta;\mathcal{D})
=
-\frac{1}{\sum_{r=1}^{N}T_r}
\sum_{r=1}^{N}
\sum_{t=1}^{T_r}
\log p_{\theta}
\left(
  w_t^{(r)}
  \mid
  w_{<t}^{(r)}
\right),
$$
ja mallin koulutus esitetään optimointiongelmana
$$
\min_{\theta} J(\theta;\mathcal{D}).
$$
Toisin kuin \(n\)-gram-mallissa, eri kontekstien
ehdolliset jakaumat eivät ole itsenäisiä parametreja,
vaan ne muodostetaan yhteisellä parametrisaatiolla
\(\theta\). Optimointiongelma ei siten hajoa erillisiksi
kontekstikohtaisiksi tehtäviksi, eikä sille yleensä ole
suljetun muodon ratkaisua.

Häviö voidaan esittää myös tokenikohtaisten tavoitejakaumien ja mallin ennustejakaumien ristientropioiden keskiarvona.
Merkitään position \(t\) tavoitejakaumaa ja mallin
ennustejakaumaa komponenteittain
$$
q_{t,j}^{(r)}
=
\mathbf{1}\{w_t^{(r)}=v_j\},
\qquad
p_{\theta,t,j}^{(r)}
=
p_{\theta}
\left(
  v_j \mid w_{<t}^{(r)}
\right),
$$
missä \(\mathbf{1}\{\cdot\}\) on indikaattorifunktio.
Koska tavoitejakauma keskittää kaiken todennäköisyysmassan
havaittuun tokeniin, saadaan
$$
-\sum_{j=1}^{|\mathcal{V}|}
q_{t,j}^{(r)}\log p_{\theta,t,j}^{(r)}
=
-\log p_{\theta}
\left(
  w_t^{(r)} \mid w_{<t}^{(r)}
\right).
$$
Negatiivisen log-uskottavuuden minimointi on täten
yhtäpitävää tavoitejakaumien ja ennustejakaumien
keskimääräisen ristientropian minimoinnin kanssa.

Koulutuksessa ehdollistetaan aineistossa havaittuihin
prefikseihin. Sekvenssiä \(w_{1:T}\) vastaava syöte
ja tavoitejono ovat
$$
x=(\langle s \rangle,w_1,\dots,w_{T}),
\qquad
y=(w_1,\dots,w_T,\langle /s \rangle).
$$
Kun \(\ell_t\) on tämän syötejonon position \(t\)
logittivektori, kausaalinen maski takaa, että
$$
\operatorname{softmax}(\ell_t)_j
=
p_{\theta}(w_t=v_j\mid w_{<t}).
$$
Kaikkien positioiden ennusteet voidaan siten laskea
rinnakkain yhdellä eteenpäinlaskennalla.

Softmax-funktion ja ristientropian yhdistelmä antaa
position häviölle yksinkertaisen derivaatan logittien
suhteen
$$
\frac{\partial}
{\partial \ell_{t,j}^{(r)}}
\left(
  -\sum_{k=1}^{|\mathcal{V}|}
  q_{t,k}^{(r)}\log p_{\theta,t,k}^{(r)}
\right)
=
p_{\theta,t,j}^{(r)}-q_{t,j}^{(r)}.
$$
Parametrien suhteen gradientti saadaan ketjusäännöllä,
ja se lasketaan vastavirta-algoritmilla, jolloin parametreja voidaan päivittää luvussa \ref{sec:grad} käsitellyllä gradienttimenetelmällä.

Laajalla tekstiaineistolla toteutettua koulutusta
kutsutaan autoregressiiviseksi esikoulutukseksi
\cite{radford2018improving}. Esikoulutettua mallia
voidaan edelleen hienosäätää esimerkiksi ohjeiden
noudattamiseen ohjevirityksellä \cite{ouyang2022training}. Valvotussa ohjevirityksessä havainto
koostuu syötekontekstista \(c\) ja tavoitevastauksesta
\(y=(y_1,\dots,y_K)\), jolloin vastaustokeneihin
rajattu negatiivinen log-uskottavuus on
$$
-\log p_{\theta}(y\mid c)
=
-\sum_{t=1}^{K}
\log p_{\theta}(y_t\mid c,y_{<t}).
$$
Optimointiperiaate säilyy samana, mutta häviö
kohdistetaan annetun kontekstin ehdollistamaan
tavoitevastaukseen.

Kun koulutettua mallia sovelletaan tekstin generointiin,
sen parametrit pidetään kiinteinä. Malli muodostaa
seuraavan tokenin todennäköisyysjakauman nykyisen
kontekstin perusteella. Jakaumasta valitaan tokeni,
joka liitetään kontekstin loppuun, ja laskenta
toistetaan näin täydentyneellä kontekstilla.

Generointi päättyy esimerkiksi lopetustokeniin tai
tuotokselle asetettuun pituusrajaan. Myös syötteen ja
tuotettujen tokenien yhteinen kontekstinpituus rajoittaa
laskentaa. Rajan saavuttamisen jälkeen jatkaminen edellyttää
kontekstin lyhentämistä tai muuta erikseen määriteltyä
kontekstinhallintaa.

Yksinkertainen valintasääntö generoitavalle tokenille on
$$
\widehat{w}_{t+1}
=
\operatorname*{arg\,max}_{v \in \mathcal{V}}
p_{\theta}(v \mid w_{1:t}),
$$
joka valitsee jokaisessa vaiheessa kyseisen jakauman
todennäköisimmän tokenin.
Vaihtoehtoisesti voidaan käyttää otantaa,
jossa seuraava tokeni arvotaan jakaumasta. Otannassa jakauman keskittyneisyyttä voidaan muuttaa
kaavalla
$$
p_{\theta,\tau}(v \mid w_{1:t})
=
\frac{\exp(\ell_{t,v}/\tau)}
{\displaystyle\sum_{u \in \mathcal{V}}
\exp(\ell_{t,u}/\tau)},
$$
missä \(\tau>0\) on lämpötilaparametri. Kun \(\tau=1\), saadaan alkuperäinen softmax-jakauma.
Pienempi lämpötila keskittää massaa suurimpien logittien
kohdalle, kun taas suurempi lämpötila tasoittaa jakaumaa.

Kielimallien sovelluksissa seuraavan tokenin ennustaminen toimii mekanismina laajempien tehtävien ratkaisemiseen. Syötekontekstissa esitetään esimerkiksi kysymys tai ohje, johon mallin odotetaan tuottavan sisällöllisesti oikea vastaus. Tällöin on erotettava toisistaan tokenijonon todennäköisyys mallin määrittämässä jakaumassa ja sen oikeellisuus annetun tehtävän ratkaisuna. Todennäköisin tokenijono ei välttämättä ratkaise tehtävää oikein, sillä autoregressiivinen koulutustavoite ei sellaisenaan takaa tuotoksen loogista tai toiminnallista oikeellisuutta.

Tämä ero korostuu luonnollisella kielellä esitettyjen kysymysten muuntamisessa SQL-kyselyiksi. Generoitu SQL-kysely voi olla syntaktisesti kelvollinen ja suoritettava, vaikka se vastaisi eri kysymykseen kuin käyttäjä tarkoitti.

\subsection{Kielimalli Text-to-SQL-tehtävissä}

Text-to-SQL-tehtävässä luonnollisella kielellä esitetty
tiedontarve muunnetaan tietokantaan kohdistuvaksi
SQL-kyselyksi. Autoregressiivisen kielimallin näkökulmasta
kyseessä on ehdollinen generointitehtävä. Tuotettavan
tokenijonon tulee muodostaa SQL-kysely, jonka merkitys
vastaa käyttäjän kysymystä käytettävissä olevan tietokannan
ja sitä koskevien määritelmien puitteissa.

Olkoon \(q\) luonnollisen kielen kysymys,
\(s\) tietokannan skeeman kuvaus ja \(m\) semanttinen
lisätieto. Skeema kuvaa esimerkiksi taulut, sarakkeet
ja niiden väliset suhteet. Semanttinen lisätieto
täydentää tätä kuvausta käsitteiden merkityksillä,
mittarien määritelmillä ja muilla tulkintasäännöillä.

Nämä tiedot esitetään mallille tokenikontekstina
$$
c=\operatorname{Serialize}(q,s,m),
$$
missä \(\operatorname{Serialize}\) tarkoittaa tietojen
muuntamista mallin käsittelemäksi tokenijonoksi.
Konteksti voi sisältää esimerkiksi tehtäväohjeen,
skeeman tekstimuotoisen kuvauksen ja esimerkkejä
kysymysten ja SQL-kyselyiden vastaavuuksista.

Kun SQL-kyselyn tokenijono on \(y=(y_1,\dots,y_K)\),
sen ehdollinen todennäköisyys voidaan kirjoittaa muodossa
$$
p_{\theta}(y \mid c)
=
\prod_{t=1}^{K}
p_{\theta}(y_t \mid c,y_{<t}).
$$
Kontekstin ja generoinnissa tarvittavan historian on
mahduttava mallin käytettävissä olevaan konteksti-ikkunaan.
Laajan tietokannan tapauksessa koko skeeman ja kaiken
lisätiedon sisällyttäminen yhteen syötteeseen ei siten
välttämättä ole tarkoituksenmukaista.

Tämä esitys kuvaa kielimallin generointivaihetta,
ei välttämättä koko Text-to-SQL-järjestelmän toimintaa.
Järjestelmä voi lisäksi valita relevantteja skeeman osia,
hakea esimerkkejä, tarkistaa tuotettua SQL-kyselyä tai
muodostaa korjatun ehdotuksen uuden mallikutsun avulla \cite{hong2024survey}.
Sovelluksen toiminta on tällöin useiden käsittelyvaiheiden
yhteistulos.

Keskeinen osatehtävä on skeemaan kohdistaminen
(engl.\ \emph{schema linking}), jossa kysymyksen ilmaukset
yhdistetään tietokannan tauluihin, sarakkeisiin ja muihin
rakenteisiin \cite{wang2020rat}. Kysymyksessä käytetty sanasto ei välttämättä
vastaa suoraan tietokannan nimiä. Lisäksi oikea vastaus
voi edellyttää useiden taulujen yhdistämistä, vaikka
käyttäjä ei mainitsisi niitä erikseen.

Tarkastellaan esimerkiksi kysymystä \emph{"Mikä oli viime vuoden myynti asiakasryhmittäin?"}
Sen muuntaminen SQL-kyselyksi edellyttää vähintään
myyntitapahtumia kuvaavan taulun, myyntimäärän
sisältävän sarakkeen, ajallisen rajauksen ja
asiakasryhmän tunnistamista. Jos asiakasryhmä sijaitsee
eri taulussa kuin myyntitapahtumat, on lisäksi
valittava oikea liitosavain ja liitospolku.
Ryhmittely ja aggregointi on muodostettava siten,
että tuloksen tarkkuustaso vastaa kysymystä.

Pelkkä skeeman tunteminen ei aina ratkaise näitä
valintoja. Esimerkin sana \emph{myynti} voi tarkoittaa
bruttomyyntiä, nettomyyntiä tai laskutettua myyntiä.
Samoin on ratkaistava, vähennetäänkö palautukset ja
perustuuko ajallinen rajaus tilauspäivään vai
laskutuspäivään. Ilmaus \emph{viime vuosi} edellyttää
viiteajankohtaa sekä tietoa siitä, tarkoitetaanko
kalenterivuotta vai tilikautta.

Skeema ja semanttinen tieto täydentävät siten toisiaan.
Skeema määrittää käytettävissä olevat rakenteet,
kun taas semanttinen tieto auttaa tulkitsemaan,
miten näitä rakenteita tulee käyttää tietyn
tiedontarpeen ratkaisemiseen. Semanttiseksi kerrokseksi
voidaan kutsua tietokannan rakenteita täydentävää
kuvausta, jossa määritellään esimerkiksi
liiketoimintakäsitteitä, mittareita ja hyväksyttyjä
taulujen välisiä suhteita \cite{CortexAnalyst}.

Kielimallille annettu semanttinen kuvaus muuttaa
generoinnin ehdollistavaa kontekstia. Se voi vähentää
monitulkintaisuutta tarjoamalla eksplisiittisiä
määritelmiä sen sijaan, että mallin olisi pääteltävä
merkitykset pelkästään sarakkeiden nimistä tai
esikoulutuksessa opituista yhteyksistä. Kuvaus ei
kuitenkaan takaa oikeaa tulosta. Malli voi soveltaa
sitä väärin, ja virheellinen tai ristiriitainen
kuvaus voi itsessään johtaa virheelliseen kyselyyn.

SQL-kyselyn oikeellisuudessa on hyödyllistä erottaa
kolme tasoa. Ensimmäinen on syntaktinen kelvollisuus:
noudattaako tuotettu kysely käytetyn SQL-murteen
kielioppia? Toinen on suoritettavuus: voidaanko kysely
suorittaa kyseisessä tietokantaympäristössä? Kolmas
on semanttinen oikeellisuus: toteuttaako kysely
käyttäjän tarkoittaman tiedonhaun?

Syntaktisesti kelvollinen kysely voi viitata
olemattomaan sarakkeeseen ja olla siksi
suorituskelvoton. Suoritettava kysely voi puolestaan
käyttää väärää aikarajausta tai liitosta ja tuottaa
siten väärän vastauksen. Lisäksi suoritettavuuteen
vaikuttavat kyselyn ulkopuoliset tekijät, kuten
käyttöoikeudet ja ympäristön tila. Suoritusvirhe
ei siten aina merkitse generoidun SQL-kyselyn
loogista virhettä.

Olkoon \(D\) tietokannan tila eli tietokantainstanssi.
Merkitään onnistuneesti suoritetun kyselyn tulosta
$$
\operatorname{Exec}(y,D).
$$
Jos \(y^{\star}\) on kysymykseen laadittu vertailukysely,
tuotettua kyselyä voidaan arvioida vertaamalla
suoritustuloksia
$$
\operatorname{Exec}(y,D)
=
\operatorname{Exec}(y^{\star},D).
$$
Vertailussa käytettävä tulosten yhtäsuuruus on
määriteltävä tehtävän mukaisesti. Esimerkiksi
rivijärjestyksen merkitys riippuu siitä,
edellyttääkö kysymys järjestettyä tulosta.

Pelkkä SQL-kyselyiden merkkijonojen yhtäsuuruus
on liian rajoittava oikeellisuuden peruste,
sillä eri tavoin kirjoitetut kyselyt voivat
toteuttaa saman tiedonhaun. Toisaalta myöskään
sama suoritustulos yhdellä tietokantainstanssilla
ei yleisesti todista kyselyiden semanttista
ekvivalenssia \cite{zhong2020semantic}. Virheellinen suodatusehto voi
esimerkiksi jäädä havaitsematta, jos aineistossa
ei ole rivejä, joihin se vaikuttaisi.

Luonnollisen kielen kysymys ei myöskään aina määritä
yksikäsitteistä tavoitekyselyä. Jos keskeinen mittari
tai aikaväli on määrittelemätön, useat SQL-kyselyt
voivat vastata kysymyksen eri tulkintoja. Tällöin
tarkentavan kysymyksen esittäminen voi olla
asianmukaisempaa kuin yhden tulkinnan valitseminen
ilman perustelua \cite{min2020ambigqa}. Jos tarvittavaa tietoa ei ole
tietokannassa, oikea toiminta voi puolestaan olla
perusteltu ilmoitus siitä, ettei kysymykseen voida
vastata käytettävissä olevilla tiedoilla.

Autoregressiivinen kielimalli tarjoaa
mekanismin SQL-kyselyiden muodostamiseen,
mutta sen tokeniennusteisiin perustuva koulutustavoite
ei yksin takaa kyselyiden oikeellisuutta \cite{ouyang2022training}.
Käytännön onnistuminen riippuu myös skeeman
esitystavasta, semanttisista määritelmistä,
kysymyksen yksikäsitteisyydestä ja järjestelmän
muista käsittelyvaiheista.

Seuraavassa luvussa määritellään tutkimusasetelma,
jossa tätä kokonaisuutta tarkastellaan Cortex Analystiin
perustuvassa Text-to-SQL-tapaustutkimuksessa.
Edellä esitetty teoria toimii tehtävän käsitteellisenä
perustana, mutta sitä ei tulkita kuvaukseksi palvelun
kaikista sisäisistä toteutusratkaisuista.
